\begin{table*}[htbp]
\centering
\caption{RQ4 pairings. $\Delta$ is domain-informed minus generic, so a positive $\Delta R^2$ favours the domain heuristic. The fourth pairing is NOT matched: random edge dropout at its configured rate keeps 69.7\% of edges against spanning forest's 58.1\%, so its gap conflates the heuristic with how much each arm removed.}
\label{tab:rq4_pairings}
\small
\resizebox{\textwidth}{!}{%
\begin{tabular}{lllrrrrr}
\toprule
\textbf{Domain-informed} & \textbf{Generic} & \textbf{Matched} & \textbf{Edge ret. (dom.)} & \textbf{Edge ret. (gen.)} & \textbf{$\Delta$ RMSE} & \textbf{$\Delta R^2$} & \textbf{$\Delta$ Spearman} \\
\midrule
Span-weighted METIS & METIS & by construction & 0.923 & 0.937 & -0.054 & 3.071 & 0.026 \\
Level slicing & Random hashing & by construction & 0.705 & 0.402 & -0.018 & 0.839 & 0.302 \\
AND-gate only & PageRank pruning & by calibration & 0.733 & 0.627 & 0.001 & -0.047 & -0.393 \\
Spanning forest & Random edge dropout & NOT MATCHED & 0.581 & 0.697 & 0.002 & -0.077 & 0.088 \\
\bottomrule
\end{tabular}
}
\\[4pt]\begin{minipage}{\linewidth}\footnotesize\emph{Each configuration is trained ONCE. These are point estimates with no run-to-run variance behind them, and RQ4's expected effect is small: a gap of a few percent cannot be distinguished from noise.}\end{minipage}
\end{table*}
